% ⚠ SUPERSEDED 2026-06-26 — FOLDED INTO docs/papers/symbolic_substrates/ (Paper 2).
% rmw3 decision "P2 absorbe a P_VSL". This skeleton's structure + abstract are the
% LANGUAGE LENS of symbolic_substrates/main.tex. Retained as source material only.
\documentclass[12pt]{article}

\usepackage[utf8]{inputenc}
\usepackage[T1]{fontenc}
\usepackage{geometry}
\geometry{a4paper, margin=2.5cm}
\usepackage{setspace}
\onehalfspacing
\usepackage{hyperref}
\hypersetup{colorlinks=true, linkcolor=black, citecolor=black, urlcolor=black}
\usepackage{natbib}
\bibliographystyle{plainnat}
\usepackage{listings}
\usepackage{xcolor}
\usepackage{booktabs}
\usepackage{amsmath}

% VSL code style
\lstdefinelanguage{VSL}{
  keywords={V, S1, S2, S3, S4, S5, viable, blocked, partial, true, false,
            recursive, phi_up, phi_down, algedonic},
  sensitive=true,
  comment=[l]{//},
  morestring=[b]",
}
\lstset{
  language=VSL,
  basicstyle=\ttfamily\small,
  keywordstyle=\bfseries,
  commentstyle=\itshape\color{gray},
  frame=single,
  breaklines=true,
}

\title{VSL: A Language for Viable Systems\\
\large and the Agent That Compiles It}

\author{[authors]}

\date{Draft — \today}

\begin{document}

\maketitle

\begin{abstract}
We describe VSL (Viable System Language), a domain-specific language derived from
Beer's Viable System Model, designed to close a gap that existing formalisms leave
open: a medium that is simultaneously parseable by a deterministic grammar, generatable
by a cognitive agent (LLM), and semantically precise enough to govern the coordination
of a multi-layer operational system. VSL is not a general-purpose notation; its grammar
encodes viability constraints directly, and its boundary conditions are determined by
the cognitive/operational layer distinction (D6) that governs the system it describes.

We argue that an LLM generating VSL is not translating but \emph{compressing}: it
extracts the operational intent from natural language and encodes it in the minimum
formal representation that preserves that intent. This agent-as-compiler framing is
distinct from classical compilation and gives rise to measurable quality criteria:
parse rate, semantic fidelity, and round-trip intent preservation. We report corpus
results (212/217, $\approx$97.7\%) and propose four controlled experiments (EXP-VSL-01
through EXP-VSL-04) to quantify compression ratio, generation fidelity, and the
Shannon entropy of VSL relative to natural language equivalents.

The closing argument connects VSL to information theory: the grammar is the codebook;
the set of valid VSL files is a low-dimensional manifold in the space of all strings;
generating VSL is finding a point on that manifold. A language is successful when it
compresses the right things — and the boundary of what VSL compresses (operational
state, not policy intent) is a design decision, not a limitation.
\end{abstract}

\textbf{Keywords:} domain-specific language; viable system model; LLM-as-compiler;
semantic compression; information theory; Shannon entropy; cybernetics; VSL

\bigskip\hrule\bigskip

% ─────────────────────────────────────────────────────────────────────────────
\section{The Problem of Saying What a System Is}
\label{sec:origin}
% ─────────────────────────────────────────────────────────────────────────────

% NARRATIVE — origin story.
% The moment we needed a new language: first handoff between sessions too long
% for a context window but too structured for pure prose.
% Why existing languages were wrong in specific ways.
% VSL emerged bottom-up from operational need, not from a language design workshop.

[SECTION TO WRITE]

% ─────────────────────────────────────────────────────────────────────────────
\section{A Grammar Derived From Structure}
\label{sec:design}
% ─────────────────────────────────────────────────────────────────────────────

% NARRATIVE + formal grammar excerpts — readable, not a manual.
% The three things VSL must express: system identity, channel contracts, viability state.
% The BNF: small enough to fit in a page. Every production rule has a reason.
% What VSL deliberately does NOT express.
% Recursion in grammar mirrors recursion in VSM.

% Example VSL fragment (seed):
\begin{lstlisting}
@vsm 0.1

V(vsf-nucleus, L1) {
  status:  viable
  channel: phi_up   -> V(vsf-s3-star, L2)
  channel: phi_down <- V(vsf-s3, L2)
  channel: algedonic -> V(vsf-s5, L3)
  recursive: false
}
\end{lstlisting}

[SECTION TO WRITE]

% ─────────────────────────────────────────────────────────────────────────────
\section{The Agent That Compiles}
\label{sec:agent}
% ─────────────────────────────────────────────────────────────────────────────

% NARRATIVE + theory — the novel contribution.
% Classical compiler vs. agent-compiler: translation vs. semantic compression.
% Measurable quality: parse rate, semantic fidelity, round-trip preservation.
% D6 as enabling constraint: keeps LLMs out of L0/L1, gives VSL clean semantics.
% Qwen as S1, Claude as S3-S5: VSL is the S2 channel between cognitive layers.

[SECTION TO WRITE]

% ─────────────────────────────────────────────────────────────────────────────
\section{Measuring the Language}
\label{sec:experiments}
% ─────────────────────────────────────────────────────────────────────────────

% EXPERIMENTAL — the evidence section.

\subsection{Existing Results}
\label{sec:existing}

% Corpus parse rate: 212/217 (~97.7%)
% --corpus sweep: BNF production coverage
% Qwen generation: vsl_run.py plan/eval/propose

[SUBSECTION TO WRITE]

\subsection{EXP-VSL-01: Compression Ratio}
\label{sec:exp01}

% NL vs. VSL token count for 10 session handoffs.
% Metric: compression_ratio = tokens_NL / tokens_VSL
% Metric: ambiguity_ratio = N_interpretations_NL / 1 (VSL by construction = 1)

[SUBSECTION TO WRITE]

\subsection{EXP-VSL-02: Round-Trip Fidelity}
\label{sec:exp02}

% intent -> Claude (NL) -> VSL -> Qwen (execution) -> outcome
% Metric: fidelity_score in [0,1]; human judge

[SUBSECTION TO WRITE]

\subsection{EXP-VSL-03: Agent Generation Quality}
\label{sec:exp03}

% 20 controlled scenarios in NL -> Qwen generates VSL
% Metric: parse_rate, semantic_accuracy, compression_ratio

[SUBSECTION TO WRITE]

\subsection{EXP-VSL-04: Shannon Entropy of VSL vs.\ NL}
\label{sec:exp04}

% H(VSL) vs H(NL) for equivalent operational content.
% Expected: H(VSL) << H(NL) — grammar constrains the space.
% D_eff_vsl: effective vocabulary dimensions.

[SUBSECTION TO WRITE]

% ─────────────────────────────────────────────────────────────────────────────
\section{What a Language Is For}
\label{sec:closing}
% ─────────────────────────────────────────────────────────────────────────────

% NARRATIVE — closing argument.
% A language is successful when it compresses the right things.
% VSL compresses operational state; S5 policy intentionally remains in NL.
% The grammar as attractor: valid VSL = low-D manifold in string space.
% Open questions: emergent language; can VSL express cost-of-existing?

[SECTION TO WRITE]

% ─────────────────────────────────────────────────────────────────────────────

\bibliography{refs}

\end{document}
